% ======================= 第 九 章 =======================
\chapter{东渡的二百一十七人}
\chaptersubtitle{ГЛАВА ДЕВЯТАЯ}{ДВЕСТИ СЕМНАДЦАТЬ}
\epigraphbox{“海中有三神山，仙人居之。”}{《史记·淮南衡山列传》}

\kaishou{本章要说明的是：徐福那条船上有二百一十七个不想再老下去的人；九百年后，他们在另一片土地上变成了神。}

\noindent
上一章末尾那个判断，落地十一年后作出。作出之后要做的事很具体：找一条船，混进去，走得越远越好。

这件事在《史记》里有一段记载，文字的着眼点完全在求仙。本章要补的是这段文字背后那两百一十七个人。

\section{徐福船队的编制}

《史记·淮南衡山列传》记徐福上书，称海中有三神山，名曰蓬莱、方丈、瀛洲，仙人居之，请得斋戒，与童男女求之。秦始皇大悦，遣振男女三千人，资之五谷种种百工而行。

这段文字有三个细节值得留意。

一是人数。三千童男女，不是一个求仙团的规模，是一个移民点的规模。

二是随行人员。五谷种子与百工，求仙用不上这两样，过日子才用得上。

三是「行」这个字。原文作「徐福得平原广泽，止王不来」，意思是找到了地方就停下来，自己称王，不回来了。

三处合起来，这份编制更像一次移民，不像一次求仙。

\section{名册中的编外人员}

`ФВ-1975-11` 是一份名册，共十四叶，记东渡随行人员，按舱位分组。

名册的前三千行是童男女，记法是编号加年龄，没有姓名。此后有二百一十七行，记法不同：每行有名字、有原职、有「自备」一栏。

这二百一十七个人不在编制之内。他们是混进去的。

名册第 12 叶的页边有一行小字，写的是「头二人自备舟具」。头二人，指排在最前的两位：一位原职记为「医药」，一位原职记为「织作」。

这两位后来各有名字：一位在日本的记载里叫八意思兼神，一位叫栉笥名。

\section{携出之物与遗留之物}

临行前，那位原职为「医药」的人带走了两样东西。

一样是天文密葬法的本身。这一样不是文书，是技法，带在身上，教不了别人，也留不下来。

一样是后羿那件事的全部档册，共十九卷。

她没有带走的是这座城的其余部分：净秽山的存货、章程总目、工籍、历法的原版。这些东西在第 VIII 章所说的那次崩溃之后被留在了原地，此后无人过问。

笔者要说清一件事：这一支没有人服过蓬莱药。他们的身体还是无秽之身，落地之后照样沾秽，照样会老。此事在《竹取物语》里被改写成了另一个样子，见下一章。

\section{口传到记纪的九百年}

抵日的时间约在前二一〇年，正当弥生中期。

《古事记》成书于七一二年，《日本书纪》成书于七二〇年，相隔八年。自抵日到《古事记》成书，中间约九百二十年。

这九百二十年里，没有文字记载。二百一十七人的后代与本地人群通婚、杂居、迁徙，口耳相传的东西一层一层沉积，到第八世纪被写成书的时候，已经凝固成了另一个样子。

\begin{table}[htbp]
\centering
\caption{三个名字的对应}
\label{tab:names}
\small
\begin{tabular}{@{}llp{62mm}@{}}
\toprule
月之都的称呼 & 日本的称呼 & 对应关系\\
\midrule
月夜见 & 月读命 & 本人未东渡，约前 1980 年死于月面；名随东渡者口传\\
八意永琳 & 八意思兼神 & 东渡二百一十七人之一，原职「医药」\\
蓬莱山辉夜 & 栉笥名·赫映姬 & 东渡二百一十七人之一，原职「织作」\\
\bottomrule
\end{tabular}
\tabsource{日本方面的名号据《古事记》与《日本书纪》本文。对应关系为本书所拟，两书原文均未言及月之都。}
\end{table}

表 \ref{tab:names} 的三行里，只有第二行是直接的：八意思兼神在记纪里主谋略、司思虑，与「医药」这个原职并不相合，但与其人在月之都的实际职掌相合。这一点第三章与第六章已有交代。

第三行最薄弱。「栉笥名」这个词在《古事记》里只出现两次，一次是名字，一次是歌谣中的双关。把它与一位织作出身的人对上，靠的是「栉笥」二字的字面，不是靠材料。

第一行根本不能算对应。月夜见没有东渡，他死在前一九八〇年的净秽山里。他的名字出现在日本的记载里，只能解释为口传：东渡的人记得这个人，讲了两百年，讲到名字与神号混在一起。

\section{可对应者只有一支}

三贵子之中，只有月读一支在记纪里可以与本书的材料对上。天照与素盏鸣没有。

这一点不奇怪。二百一十七人进入的是一个有几千年历史的人群，他们的后代在九百年后只占极小的一部分。能被记住的，只有那些被反复讲述的名字；讲述的人多了，名字就神格化；讲述的人少了，名字就散掉。

三贵子里留下了一支，说明这一支被讲得最多。至于为什么被讲得最多，材料里没有答案。

\section{本章结论}

本章说了四件事。

徐福那条船是一次移民，不是一次求仙。三千童男女、五谷、百工，是过日子的编制。

船上有二百一十七个不在编制里的人，领头的是两位，原职分别为医药与织作。

他们带走了技法与档册，留下了城的其余部分。

九百年后，三个人名凝固成三个神号。

本章的自我拆台：二百一十七这个数，只有一份名册做主，而那份名册的来历本身不清楚。它是如何进入 `ФВ` 卷宗的，卷宗里没有交代。若这份名册是后人伪作，本章的整个第二部分就得推翻。

笔者保留它，理由与第三章相同：它与另外两处材料对得上。名册所记的原职分布，与第 I 卷简二十二所记的工种分类一致；名册所记的舱位分组，与第八章玉山十二舱的编号方式一致。伪造者要同时对上这两处，成本太高。
